\documentclass[../main.tex]{subfiles}
\begin{document}

\section{Overview}


\section{Datasets}
\label{sec:datasets}
\subsection{Synthetic analogies from BATS}
The Bigger Analogy Test Set (BATS) is a benchmark for probing the relational and semantic
knowledge of word embeddings and language models. We extract the entities in its
lexicography category to build a synthetic dataset of ten prompts per subcategory. Each
prompt has the form ``The saying goes: $A$ is to $B$ as $C$ is to $D$'', where $A$, $B$,
$C$, $D$ are entities of one subcategory, and the task is to predict $D$ from $A$, $B$,
and $C$.

\subsection{Multi-hop reasoning}
We follow the multi-hop queries of Ameisen et al.~\cite{ameisen2025circuit}. Each query
requires reasoning across several steps to reach the answer and so tests the model's
ability to chain basic inferences. In both datasets we keep only prompts the model answers
correctly with confidence above $20\%$, so every attribution graph explains a prediction
the model actually makes.

\section{Experimental Setup}
\label{sec:setup}
\subsection{Environment}
Experiments run on a single workstation with an NVIDIA GeForce RTX 4090 GPU (24\,GB
VRAM) and 64\,GB of RAM, under Ubuntu 22.04 LTS. All code uses Python 3.10.

\subsection{Tools and libraries}
We generate the attribution graphs with the circuit-tracing library of Ameisen et
al.~\cite{ameisen2025circuit}, reuse its feature metadata, and adapt its viewer to
visualize the summary graphs.

\section{Evaluation Methodology}
\label{sec:eval-method}
\subsection{Pruning evaluation}
We evaluate the pruning stage along two axes: faithfulness to the original attribution
graph and interpretability of the retained nodes. Faithfulness is measured by graph
completeness and replacement score after pruning. Graph completeness measures how much of
the original attribution-graph computation remains explainable by the retained subgraph,
while replacement score measures whether the pruned graph still preserves the model
prediction under the replacement-model evaluation. Interpretability is measured by the
number and proportion of retained syntactic features, because these nodes usually explain
prompt scaffolding, punctuation, or function-word structure rather than the semantic
relation needed for the analogy task.

Figures~\ref{fig:prune-alpha-entmax-no-token-attr} and
\ref{fig:prune-alpha-softmax-no-token-attr} compare graph completeness and replacement
score against the average number of retained nodes. The curves separate the mixture
parameter $\alpha$, where $\alpha=0$ corresponds to influence-only pruning and larger
values add SHAP-based relevance to the node-importance score. The result shows that adding
relevance does not substantially reduce graph-level faithfulness at matched graph sizes.
The effect is most useful in the small-graph regime, where the pruned graph must keep only
the nodes that are most likely to support the target prediction.

\begin{figure}[ht]
\centering
\includegraphics[width=0.95\linewidth]{Figure/prune_alpha_sweep_entmax_no_token_attr.png}
\caption{Pruning alpha sweep with entmax normalization. The figure reports graph
completeness and replacement score against average retained node count, omitting token
attribution faithfulness to focus on graph-level faithfulness.}
\label{fig:prune-alpha-entmax-no-token-attr}
\end{figure}

\begin{figure}[ht]
\centering
\includegraphics[width=0.95\linewidth]{Figure/prune_alpha_sweep_softmax_no_token_attr.png}
\caption{Pruning alpha sweep with softmax normalization. The figure reports graph
completeness and replacement score against average retained node count, omitting token
attribution faithfulness to focus on graph-level faithfulness.}
\label{fig:prune-alpha-softmax-no-token-attr}
\end{figure}

Table~\ref{tab:syntactic-feature-count-template} defines the manual counting table used
to compare influence-only pruning against ours setting $\alpha=0.5$. A retained
feature is counted as syntactic when its label primarily refers to punctuation, stop words,
inflection, word order, prompt-format tokens, or other grammatical scaffolding. It is not
counted as syntactic when its label refers to the semantic relation, the source entity, the
target entity, or the predicted answer. The final count therefore tests whether our score removes graph mass that is faithful but not useful for a compact semantic
explanation.

\begin{table}[ht]
\centering
\caption{Template for counting syntactic features after pruning. Counts are computed over
the retained feature nodes of the analogy graphs at matched graph-size settings.}
\label{tab:syntactic-feature-count-template}
\begin{tabular}{llrrrrr}
\toprule
Pruning score & Normalization & $\alpha$ & Stop-word & Punctuation & Other syntactic & Total syntactic \\
\midrule
Influence only & Entmax 1.25 & 0.0 & -- & -- & -- & -- \\
Ours alpha 0.5 & Entmax 1.25 & 0.5 & -- & -- & -- & -- \\
Influence only & Softmax & 0.0 & -- & -- & -- & -- \\
Ours alpha 0.5 & Softmax & 0.5 & -- & -- & -- & -- \\
\bottomrule
\end{tabular}
\end{table}
\subsection{Clustering evaluation}
We evaluate clustering on the entmax-pruned analogy graphs by comparing the saved
summary-graph partition against four matched-size baselines. Each baseline uses the same
pruned graph and the same number of feature supernodes as the saved summary graph. The
silhouette score measures role-vector compactness, while internalized edge mass measures
the fraction of pruned edge mass hidden inside supernodes. DAG loss measures the
additional superedge mass removed when the summary adjacency is converted into an
acyclic graph.

\begin{table}[ht]
\centering
\caption{Clustering evaluation on 100 entmax-pruned analogy graphs. Lower internalized
edge mass and lower DAG loss indicate better preservation of visible causal structure.}
\label{tab:clustering-baseline-entmax}
\begin{tabular}{lrrr}
\toprule
Solver & Silhouette score (raw) & Internalized edge mass & DAG loss \\
\midrule
Ours & 0.0409 & 0.0491 & 0.0984 \\
Spectral on role vectors & -0.0347 & 0.1081 & 0.1257 \\
K-means on role vectors & 0.5058 & 0.0862 & 0.1423 \\
Spectral on adjacency & -0.4263 & 0.1590 & 0.1818 \\
Random same-size clusters & -0.5908 & 0.0762 & 0.2058 \\
\bottomrule
\end{tabular}
\end{table}

Table~\ref{tab:clustering-baseline-entmax} shows that K-means produces the most compact
role-vector clusters, but this compactness hides more causal structure and loses more
mass during DAG construction. The proposed clustering has a lower raw silhouette score
than K-means, but it keeps the lowest internalized edge mass and the lowest DAG loss.
This result is consistent with the objective of the summary graph: the clustering should
not merely maximize geometric compactness, but should preserve causal information in the
visible supernode graph.

\subsection{Labels evaluation}
\subsection{Steering case studies}
A summary graph is faithful only if its supernode edges reflect genuine causal
dependencies. We test this directly with a steering intervention: we perturb one
supernode at its source and check whether every other supernode moves in the direction
its incoming edge predicts. The test grades the clustering partition together with the
edges $W_H$ at once, so it scores the summary graph as a whole rather than any single
stage.

\paragraph{Intervention protocol.}
For each non-fixed supernode $C$ we negate every member feature, setting its activation
to $M \cdot a$ with steer factor $M = -1$ and $a$ the original activation, then run one
forward pass with all members of $C$ steered together. The perturbation is applied as
\emph{constrained} (direct-effect) patching, following Ameisen et
al.~\cite{ameisen2025circuit}: each steered feature writes its decoder vector only over
the layer window $[l, l+1]$ around its own layer $l$ --- its own-layer reconstruction
plus the first cross-layer write --- while the attention patterns and LayerNorm scales
are frozen at their clean values. Freezing attention and normalization isolates the
\emph{direct} effect that the summary edges model and suppresses the indirect paths the
attribution graph does not represent.

\paragraph{Measured response.}
After the intervention we read the activation change $\Delta a_{C'}$ of every other
feature supernode $C'$ and the probability change $\Delta p_{C'}$ of every logit
supernode and of the target token. For a steered source $C$ and a downstream supernode
$C'$ joined by the summary edge $W_{H,\,C \to C'}$, the edge \emph{predicts} that the
downstream response carries the sign
$\operatorname{sgn}(W_{H,\,C \to C'}) \cdot \operatorname{sgn}(\Delta a_C)$, where
$\Delta a_C$ is the activation change of the steered supernode itself (negative under
negation).


\section{Chapter Summary}


\end{document}
